\documentclass[11pt,oneside]{article}
\usepackage[a4paper,margin=27mm,headheight=15pt]{geometry}
\usepackage[T1]{fontenc}
\usepackage[utf8]{inputenc}
\IfFileExists{libertinus.sty}{\usepackage{libertinus}}{\usepackage{lmodern}}
\usepackage{microtype}
\usepackage{amsmath,amssymb,amsthm,mathtools,mathrsfs}
\usepackage{bm}
\usepackage{booktabs,tabularx,array}
\usepackage{xcolor}
\usepackage{enumitem}
\usepackage{fancyhdr}
\usepackage{titlesec}
\usepackage{xurl}
\usepackage{hyperref}
\usepackage[nameinlink,noabbrev]{cleveref}

\definecolor{linkblue}{RGB}{25,75,135}
\definecolor{shadegray}{RGB}{246,247,249}
\definecolor{rulegray}{RGB}{195,201,208}
\hypersetup{
  colorlinks=true,
  linkcolor=linkblue,
  citecolor=linkblue,
  urlcolor=linkblue,
  pdftitle={Decay of the Fourier Transform of Rvachev's Up-Function},
  pdfsubject={Dyadic products, ergodic sums, Thue--Morse polynomials, and sharp decay envelopes},
  pdfkeywords={Rvachev up-function, Fabius function, Fourier transform, sinc product, Thue-Morse, Gelfond exponent, lacunary product, ergodic theorem}
}
\urlstyle{same}

\pagestyle{fancy}
\fancyhf{}
\fancyhead[L]{\small Fourier decay of Rvachev's up-function}
\fancyhead[R]{\small\nouppercase{\leftmark}}
\fancyfoot[C]{\thepage}
% Canonical project-wide mathematical notation.
% Canonical mathematical notation for active FabiusFunction LaTeX documents.
%
% This file is intentionally a plain TeX input rather than a LaTeX package.
% Every standalone document loads it by an explicit path relative to that
% document's directory; included fragments inherit the definitions from their
% root document.  Keep this file independent of document layout and metadata.
%
% Contract:
%   * one command has one mathematical meaning throughout the active corpus;
%   * names encode semantic roles rather than a document-local choice of letter;
%   * visually similar but differently normalized objects have different print;
%   * every command below is defined normatively in the notation catalogue.
%
% The active corpus excludes docs/archive/.  Historical sources may retain
% their original notation only inside an explicitly labelled quotation.

\makeatletter
\@ifundefined{mathbb}{%
  \PackageError{fabius-notation}{The amssymb package must be loaded first}%
    {Load amsmath and amssymb before inputting fabius-notation.tex.}%
}{}
\@ifundefined{operatorname}{%
  \PackageError{fabius-notation}{The amsmath package must be loaded first}%
    {Load amsmath before inputting fabius-notation.tex.}%
}{}
\makeatother

% Version stamp used by the catalogue and by static audits.
\newcommand{\FabiusNotationVersion}{2026-08-31}

% Number systems.  NaturalNumbers is explicitly the nonnegative convention.
\newcommand{\NaturalNumbers}{\mathbb{N}_{0}}
\newcommand{\PositiveIntegers}{\mathbb{N}_{>0}}
\newcommand{\IntegerNumbers}{\mathbb{Z}}
\newcommand{\RationalNumbers}{\mathbb{Q}}
\newcommand{\RealNumbers}{\mathbb{R}}
\newcommand{\ComplexNumbers}{\mathbb{C}}
\newcommand{\ExtendedRealNumbers}{\overline{\mathbb{R}}}

% The three principal Fabius--Rvachev functions.  Subscripts are deliberate:
% a bare F and a calligraphic F have several unrelated uses in the corpus.
\newcommand{\FabiusBounded}{F_{\mathrm{bd}}}
\newcommand{\FabiusGlobal}{F_{\mathrm{gl}}}
\newcommand{\FabiusQuantile}{Q_{F}}
\newcommand{\FabiusClampedQuantile}{Q_{F,\mathrm{cl}}}
\newcommand{\RvachevUp}{\operatorname{up}}

% Constants, elementary operators, and delimiters.
\newcommand{\EulerE}{\mathrm{e}}
\newcommand{\ImaginaryUnit}{\mathrm{i}}
\newcommand{\Differential}{\mathop{}\!\mathrm{d}}
\newcommand{\AbsoluteValue}[1]{\left\lvert #1\right\rvert}
\newcommand{\Norm}[1]{\left\lVert #1\right\rVert}
\newcommand{\Floor}[1]{\left\lfloor #1\right\rfloor}
\newcommand{\Ceiling}[1]{\left\lceil #1\right\rceil}
\newcommand{\FractionalPart}[1]{\left\{#1\right\}}
\newcommand{\PositivePart}[1]{\left(#1\right)_{+}}
\newcommand{\IndicatorOf}[1]{\mathbf{1}_{#1}}
\newcommand{\SetOf}[1]{\left\{#1\right\}}
\newcommand{\InnerProduct}[1]{\left\langle #1\right\rangle}
\newcommand{\CardinalityOf}[1]{\left\lvert #1\right\rvert}
\newcommand{\DefinitionEquals}{\mathrel{\coloneqq}}
\newcommand{\Sign}{\operatorname{sgn}}
\newcommand{\IdentityOperator}{\operatorname{id}}
\newcommand{\SpanOperator}{\operatorname{span}}
\newcommand{\TraceOperator}{\operatorname{tr}}
\newcommand{\RankOperator}{\operatorname{rank}}
\newcommand{\DiagonalOperator}{\operatorname{diag}}
\newcommand{\ResidueOperator}{\operatorname*{Res}}
\newcommand{\LeastCommonMultiple}{\operatorname{lcm}}
\newcommand{\OrderOperator}{\operatorname{ord}}
\newcommand{\PrincipalLogarithm}{\operatorname{Log}}
\newcommand{\FallingFactorial}[2]{\left(#1\right)^{\underline{#2}}}
\newcommand{\RisingFactorial}[2]{\left(#1\right)^{\overline{#2}}}

% Support, topology, convolution, and metric notation.
\newcommand{\SupportOperator}{\operatorname{supp}}
\newcommand{\SupportOf}[1]{\SupportOperator\!\left(#1\right)}
\newcommand{\TopologicalSupportOf}[1]{\operatorname{tsupp}\!\left(#1\right)}
\newcommand{\Convolution}{\mathbin{*}}
\newcommand{\MetricDistance}{\operatorname{dist}}
\newcommand{\TotalVariationDistance}{d_{\mathrm{TV}}}

% Probability.  Equality and convergence in law are not metric distance.
\newcommand{\Probability}{\mathbb{P}}
\newcommand{\Expectation}{\mathbb{E}}
\newcommand{\Variance}{\operatorname{Var}}
\newcommand{\Covariance}{\operatorname{Cov}}
\newcommand{\LawOperator}{\operatorname{Law}}
\newcommand{\LawOf}[1]{\LawOperator\!\left(#1\right)}
\newcommand{\EqualInLaw}{\mathrel{\overset{\mathrm{law}}{=}}}
\newcommand{\ConvergesInLaw}{\mathrel{\xrightarrow{\mathrm{law}}}}

% Fourier and integral transforms.  The printed labels expose normalization.
% FourierTwoPi uses exp(-2*pi*i*x*xi); FourierAngular uses exp(-i*x*omega).
\newcommand{\FourierTwoPi}[1]{\widehat{#1}_{\,2\pi}}
\newcommand{\FourierAngular}[1]{\widehat{#1}_{\,\mathrm{ang}}}
\newcommand{\LaplaceTransformSymbol}{\mathcal{L}}
\newcommand{\MellinTransformSymbol}{\mathcal{M}}
\newcommand{\LaplaceTransformOf}[1]{\LaplaceTransformSymbol\!\left[#1\right]}
\newcommand{\MellinTransformOf}[1]{\MellinTransformSymbol\!\left[#1\right]}
\newcommand{\MomentGeneratingFunction}[1]{M_{#1}}
\newcommand{\CharacteristicFunction}[1]{\varphi_{#1}}

% The two standard sinc functions must never share one symbol.
\newcommand{\SincRad}{\operatorname{sinc}_{\mathrm{rad}}}
\newcommand{\SincPi}{\operatorname{sinc}_{\pi}}

% Binary arithmetic and Thue--Morse notation.
\newcommand{\BinaryDigitSum}{\operatorname{wt}_{2}}
\newcommand{\ThueMorseSignSymbol}{\varepsilon}
\newcommand{\ThueMorseSign}[1]{\varepsilon_{#1}}
\newcommand{\PadicValuation}[1]{\nu_{#1}}
\newcommand{\TwoAdicValuation}{\nu_{2}}

% q-series notation.  Every base is an explicit argument.
\newcommand{\QPochhammer}[3]{\left(#1;#2\right)_{#3}}
\newcommand{\GaussianBinomial}[3]{%
  \genfrac{[}{]}{0pt}{}{#1}{#2}_{#3}%
}
\newcommand{\QInteger}[2]{\left[#1\right]_{#2}}

% Named special functions.
\newcommand{\Polylogarithm}{\operatorname{Li}}
\newcommand{\LambertW}{W}
\newcommand{\GammaFunction}{\Gamma}
\newcommand{\RiemannZeta}{\zeta}

% Asymptotics.  BigO and LittleO are operator tokens for existing formula
% grammar; prefer the At forms so that the limiting regime is printed.
\newcommand{\BigO}{\mathcal{O}}
\newcommand{\LittleO}{o}
\newcommand{\BigOAt}[2]{\mathcal{O}_{#2}\!\left(#1\right)}
\newcommand{\LittleOAt}[2]{o_{#2}\!\left(#1\right)}

\renewcommand{\headrulewidth}{0.4pt}

\titleformat{\section}{\Large\bfseries}{\thesection}{0.7em}{}
\titleformat{\subsection}{\large\bfseries}{\thesubsection}{0.7em}{}
\titleformat{\subsubsection}{\normalsize\bfseries}{\thesubsubsection}{0.7em}{}
\setcounter{secnumdepth}{3}
\setcounter{tocdepth}{2}
\setlist[itemize]{leftmargin=2em,itemsep=0.25em,topsep=0.4em}
\setlist[enumerate]{leftmargin=2.3em,itemsep=0.25em,topsep=0.4em}
\setlength{\emergencystretch}{2em}
\allowdisplaybreaks

\newtheorem{theorem}{Theorem}[section]
\newtheorem{proposition}[theorem]{Proposition}
\newtheorem{lemma}[theorem]{Lemma}
\newtheorem{corollary}[theorem]{Corollary}
\theoremstyle{definition}
\newtheorem{definition}[theorem]{Definition}
\newtheorem{example}[theorem]{Example}
\theoremstyle{remark}
\newtheorem{remark}[theorem]{Remark}
\newtheorem{warning}[theorem]{Warning}

\newcommand{\dist}{\mathop{\longrightarrow}\limits^{\mathrm d}}
\newcommand{\T}{T}

\newenvironment{boxedremark}{%
  \par\smallskip\noindent\begin{center}\begin{minipage}{0.94\textwidth}%
  \color{black}\hrule\vspace{0.6em}\small
}{%
  \vspace{0.6em}\hrule\end{minipage}\end{center}\smallskip
}

\title{\bfseries Decay of the Fourier Transform of Rvachev's Up-Function\\[0.35em]
\large Dyadic products, ergodic fluctuations, and sharp annular envelopes}
\author{A rigorous analysis for Vladimir Reshetnikov's \texttt{ProveIt} project}
\date{26 August 2026}

\begin{document}
\maketitle

\begin{abstract}
Let
\[
 f(x)=\widehat{\RvachevUp}(x)
 =\prod_{n=0}^{\infty}\frac{\sin(\pi x/2^n)}{\pi x/2^n}
\]
be the Fourier transform of Rvachev's compactly supported up-function, with the
Fourier convention $\widehat{\RvachevUp}(x)=\int_{\RealNumbers}\RvachevUp(t)\EulerE^{-2\pi\ImaginaryUnit xt}\Differential t$.
The phrase ``the decay rate of $f$'' is intrinsically ambiguous: $f$ vanishes at
every nonzero integer, its values on a fixed dyadic ray are governed by a
Birkhoff sum for the doubling map, and the largest lobe in a dyadic annulus is
an extremal problem for a Thue--Morse trigonometric polynomial.  These regimes
have different second-order behavior.

Writing $a=\log 2$, we prove the exact dyadic decomposition that reduces the
problem to the potential $\phi(t)=\log|\sin\pi t|$.  For Lebesgue-almost every
mantissa $y\in(1,2)$,
\[
 \log|f(2^Ny)|
 =-\frac{\bigl(\log(2^Ny)\bigr)^2}{2a}
  -\left(\frac{\log\pi}{a}+\frac32\right)\log(2^Ny)
  +\LittleO(N).
\]
The centered remainder satisfies a central limit theorem with variance
$\pi^2/4$ per dyadic step and a law of the iterated logarithm.  Thus even the
natural deterministic normalization does not converge to a nonzero constant.
The slowest invariant average is attained by the period-two orbit
$\{1/3,2/3\}$, producing the power exponent
\[
 \kappa_*=\frac{\log\pi+a/2-\log(\sqrt3/2)}{a}
 =\frac{\log\pi}{\log2}+\frac32-\frac{\log3}{2\log2}.
\]
Finally, if
$M_N=\max_{2^N\le x\le2^{N+1}}|f(x)|$, then the classical Gelfond estimate and
a boundary-layer optimization give
\[
 \begin{split}
 \log M_N={}&-\frac a2N^2
 +\bigl(\log(\sqrt3/2)-\log\pi-a/2\bigr)N\\
 &-\frac{(\log N)^2}{2a}
 +\frac{\log N\,\log\log N}{a}
 +\BigO(\log N).
 \end{split}
\]
Equivalently, the annular envelope has a genuine
$-(\log\log x)^2/(2\log2)$ correction, followed by a
$(\log\log x)(\log\log\log x)/\log2$ term.  This rules out both a global
pointwise equivalent and the fixed logarithmic-power correction proposed in
the starting draft.
\end{abstract}

\tableofcontents
\newpage

\section{The question and the form of the answer}

Rvachev's up-function $\RvachevUp$ is an even, nonnegative $C^\infty$ function
supported on $[-1,1]$, normalized by $\int_{\RealNumbers}\RvachevUp=1$.  It is the density of an
infinite sum of independent centered uniform random variables, and its Fourier
transform is the infinite sinc product studied below; see
\cite{arias1982,proveit,wikipedia}.  Smooth compact support already implies
that $f=\widehat{\RvachevUp}$ decreases faster than every inverse power on the real
axis.  The product contains much more information: its dominant logarithmic
scale is quadratic in $\log x$.

A preliminary draft in the \texttt{ProveIt} repository proposed a single
pointwise equivalent of the form
\begin{equation}
 |f(x)|\stackrel{?}{\sim}
 C(\log x)^{-p}x^{-\log(2\pi)/\log2}
 \exp\!\left(-\frac{(\log x)^2}{2\log2}\right).
 \label{eq:draft-form}
\end{equation}
That form cannot be correct as a pointwise statement.  The first obstruction is
immediate: $f(m)=0$ for every nonzero integer $m$.  The more substantive
obstruction is dynamical.  On writing $x=2^Ny$ with $1\le y<2$, the logarithm
of the product contains
\[
 \sum_{k=1}^{N}\log|\sin(\pi2^k y)|,
\]
a lacunary Birkhoff sum.  It has a typical mean, Gaussian-scale fluctuations,
exceptional periodic averages, and an extremal average.  No deterministic
normalizer can erase all of these simultaneously.

There are therefore several mathematically natural questions.
\begin{enumerate}
\item What is the decay along a fixed dyadic ray $x_N=2^Ny$?
\item What happens for Lebesgue-typical mantissas $y$?
\item Which fixed rays decay most slowly?
\item How large can any lobe be in the dyadic annulus $[2^N,2^{N+1}]$?
\end{enumerate}
The article answers all four.  The last question is the appropriate
interpretation of a global upper envelope, while the second describes the
amplitude seen at a typical point of a large dyadic annulus.

\begin{boxedremark}
\textbf{Main constants.}
Throughout, put
\[
 a=\log2,\qquad
 \beta=\log\frac{\sqrt3}{2},
\]
and define
\begin{align}
 \kappa_{\mathrm{typ}}
 &=\frac{\log\pi+3a/2}{a}
   =\frac{\log\pi}{\log2}+\frac32
   \approx3.1514961295,                                      \label{eq:ktyp}\\
 \kappa_*
 &=\frac{\log\pi+a/2-\beta}{a}
   =\frac{\log\pi}{\log2}+\frac32-
     \frac{\log3}{2\log2}
   \approx2.3590148791.                                      \label{eq:kstar}
\end{align}
The first is the typical power exponent.  The second comes from the maximal
Birkhoff average and controls the linear-in-$\log x$ term of the slow annular
envelope.
\end{boxedremark}

\section{Product structure, zeros, signs, and lobes}

We use the entire extension
\begin{equation}
 f(z)=\prod_{n=0}^{\infty}\SincRad\!\left(\frac{\pi z}{2^n}\right),
 \qquad \SincRad w=\begin{cases}\sin w/w,&w\ne0,\\1,&w=0.\end{cases}
 \label{eq:f-product}
\end{equation}
The product converges locally uniformly because
$\log\SincRad w=-w^2/6+\BigO(w^4)$ near the origin.

\begin{proposition}[Equivalent products and dilation equations]
For every $z\in\ComplexNumbers$,
\begin{align}
 f(z)
 &=\prod_{m=1}^{\infty}
   \left(\cos\frac{\pi z}{2^m}\right)^m,                    \label{eq:cos-product}\\
 &=\prod_{m=1}^{\infty}
   \left(1-\frac{z^2}{m^2}\right)^{1+\TwoAdicValuation(m)}.             \label{eq:weier-product}
\end{align}
Moreover,
\begin{align}
 f(z)&=\SincRad(\pi z)f(z/2),                                  \label{eq:feq-down}\\
 f(2z)&=\SincRad(2\pi z)f(z).                                  \label{eq:feq-up}
\end{align}
\end{proposition}

\begin{proof}
The identity $\SincRad w=\prod_{j\ge1}\cos(w/2^j)$ follows by iterating the
double-angle formula and passing to the limit.  In the product over $n\ge0$,
a factor $\cos(\pi z/2^m)$ occurs once for every decomposition $m=n+j$ with
$n\ge0$ and $j\ge1$, hence exactly $m$ times.  This proves
\eqref{eq:cos-product}.

Euler's sine product gives
\[
 \SincRad\!\left(\frac{\pi z}{2^n}\right)
 =\prod_{k=1}^{\infty}\left(1-\frac{z^2}{(2^nk)^2}\right).
\]
A positive integer $m$ can be written as $m=2^nk$ for precisely
$1+\TwoAdicValuation(m)$ values of $n\ge0$.  Regrouping the locally uniformly convergent
product proves \eqref{eq:weier-product}.  Finally, separating the first factor
in \eqref{eq:f-product}, or replacing $z$ by $2z$, gives
\eqref{eq:feq-down}--\eqref{eq:feq-up}.
\end{proof}

\begin{proposition}[Zeros and Thue--Morse signs]
The nonzero zeros of $f$ are exactly the integers.  The zero at
$m\in\IntegerNumbers\setminus\{0\}$ has multiplicity $1+\TwoAdicValuation(|m|)$.  If $x>0$ is not an
integer and $n=\lfloor x\rfloor$, then
\begin{equation}
 \Sign f(x)=(-1)^{\BinaryDigitSum(n)},                                   \label{eq:sign-tm}
\end{equation}
where $\BinaryDigitSum(n)$ is the sum of the binary digits of $n$.
\end{proposition}

\begin{proof}
The zero statement follows immediately from \eqref{eq:weier-product}.  Starting
with $f(0)=1$, the sign on $(n,n+1)$ is changed by the total multiplicity
\[
 \sum_{m=1}^{n}(1+\TwoAdicValuation(m))
 =n+\TwoAdicValuation(n!)=n+(n-\BinaryDigitSum(n))=2n-\BinaryDigitSum(n).
\]
Its parity is the parity of $\BinaryDigitSum(n)$, proving \eqref{eq:sign-tm}.
\end{proof}

The zeros divide the graph into unit lobes.  Every lobe after the initial
one has exactly one critical point; the initial lobe decreases from its endpoint
maximum at the origin.  This is useful when an ``envelope'' is interpreted as
the sequence of lobe maxima.

\begin{theorem}[Exactly one extremum in each noninitial unit interval]
For each integer $n\ge1$, the function $|f|$ has exactly one critical point in
$(n,n+1)$.  It is a strict maximum.  Equivalently, the equation
\begin{equation}
 \sum_{m=1}^{\infty}\frac{1+\TwoAdicValuation(m)}{m^2-x^2}=0            \label{eq:critical-equation}
\end{equation}
has exactly one solution in $(n,n+1)$ for every $n\ge1$.
\end{theorem}

\begin{proof}
Away from the integers, logarithmic differentiation of
\eqref{eq:weier-product} gives
\[
 \frac{f'(x)}{f(x)}=-2xH(x),\qquad
 H(x)=\sum_{m=1}^{\infty}\frac{1+\TwoAdicValuation(m)}{m^2-x^2}.
\]
The series and its derivative converge locally uniformly off the integers,
because $\TwoAdicValuation(m)=\BigO(\log m)$.  For $x>0$,
\[
 H'(x)=2x\sum_{m=1}^{\infty}
 \frac{1+\TwoAdicValuation(m)}{(m^2-x^2)^2}>0.
\]
For $n\ge1$, on $(n,n+1)$ one has $H(x)\to-\infty$ as $x\downarrow n$ and
$H(x)\to+\infty$ as $x\uparrow n+1$.  Hence $H$ has exactly one zero there.
At that zero the derivative of $f'/f$ is negative, so $\log|f|$, and therefore
$|f|$, has a strict local maximum.  On $(0,1)$ every denominator in $H$ is
positive, so $f'/f<0$ there; the initial lobe decreases strictly from $f(0)=1$
to $f(1)=0$.
\end{proof}

\begin{remark}
Let $A_n=\max_{n\le x\le n+1}|f(x)|$.  Then the annular maximum introduced
later is
\[
 M_N=\max_{2^N\le n<2^{N+1}}A_n.
\]
Thus the annular theorem is also a theorem about the upper envelope of the
unique extrema in successive lobes.
\end{remark}

\section{The exact dyadic decomposition}

The decisive step is not an approximation.  It is an exact decomposition into
a deterministic quadratic term and a dynamical sum.

Let
\begin{equation}
 \T(t)=2t\pmod1,
 \qquad
 \phi(t)=\log|\sin\pi t|,
 \qquad
 \psi(t)=\log|2\sin\pi t|=\phi(t)+a.                       \label{eq:potentials}
\end{equation}
Both potentials are understood as $-\infty$ at $t=0$; they belong to $L^p(0,1)$
for every finite $p$.

\begin{theorem}[Exact dyadic identity]
Let $N\ge0$, $1\le y<2$, and $z=y-1$.  Then
\begin{equation}
 \boxed{\quad
 \log|f(2^Ny)|
 =\log|f(y)|-N\log(\pi y)-\frac{aN(N+1)}2
  +\sum_{k=1}^{N}\phi(\T^kz).
 \quad}                                                     \label{eq:exact-dyadic}
\end{equation}
The equality is interpreted as $-\infty=-\infty$ when an orbit point is zero.
\end{theorem}

\begin{proof}
Split the product \eqref{eq:f-product} after the $N$ factors whose arguments
are at least of order one:
\[
 f(2^Ny)=f(y)\prod_{k=1}^{N}\SincRad(\pi2^ky).
\]
Because $2^k$ is an integer,
$|\sin(\pi2^ky)|=|\sin(\pi2^kz)|$.  Taking absolute values and logarithms gives
\eqref{eq:exact-dyadic}, since $\sum_{k=1}^{N}k=N(N+1)/2$.
\end{proof}

\begin{corollary}[Uniform log-normal upper scale]
As $x\to+\infty$,
\begin{equation}
 \log|f(x)|\le-\frac{(\log x)^2}{2\log2}+\BigO(\log x),     \label{eq:universal-upper}
\end{equation}
with an absolute implied constant.  Consequently, for every $A>0$,
$f(x)=\BigO_A(x^{-A})$.
\end{corollary}

\begin{proof}
Write $x=2^Ny$ with $1\le y<2$ and use
$|f(y)|\le1$ and $\phi\le0$ in \eqref{eq:exact-dyadic}.  Since
$N=(\log x-\log y)/a$, the quadratic term is
$-(\log x)^2/(2a)+\BigO(\log x)$.  The last assertion follows because a
negative quadratic in $\log x$ eventually dominates $-A\log x$.
\end{proof}

\begin{warning}
Equation \eqref{eq:universal-upper} is a uniform upper bound, not a uniform
equivalent.  There can be exceptionally close returns of $\T^kz$ to zero, and
at dyadic rational points the product vanishes exactly.  A matching lower bound
cannot hold pointwise on the whole real axis.
\end{warning}

\section{Fixed dyadic rays and their Birkhoff averages}

For $y\in(1,2)$ put
\[
 z=y-1,\qquad w=\T z,
\]
so that
\[
 \sum_{k=1}^{N}\phi(\T^kz)
 =\sum_{j=0}^{N-1}\phi(\T^jw).
\]
This shift makes the standard Birkhoff notation convenient.

\begin{theorem}[Ray law from a finite Birkhoff average]
Suppose $f(y)\ne0$ and that the finite limit
\begin{equation}
 \alpha(y)=\lim_{N\to\infty}\frac1N
 \sum_{j=0}^{N-1}\phi(\T^jw)                               \label{eq:alpha-def}
\end{equation}
exists.  For $x_N=2^Ny$,
\begin{equation}
 \log|f(x_N)|
 =-\frac{(\log x_N)^2}{2a}
  -\kappa(\alpha(y))\log x_N+\LittleO(\log x_N),            \label{eq:ray-law}
\end{equation}
where
\begin{equation}
 \kappa(\alpha)=\frac{\log\pi+a/2-\alpha}{a}.              \label{eq:kappa-alpha}
\end{equation}
\end{theorem}

\begin{proof}
By \eqref{eq:exact-dyadic},
\[
 \log|f(2^Ny)|
 =-\frac a2N^2+
 N\bigl(\alpha(y)-\log(\pi y)-a/2\bigr)+\LittleO(N).
\]
Substitute $N=(\log x_N-\log y)/a$.  The terms involving
$\log y\cdot\log x_N$ cancel, leaving precisely
\eqref{eq:ray-law}--\eqref{eq:kappa-alpha}.
\end{proof}

The cancellation of $\log y$ is important: once the result is expressed in
terms of the actual point $x_N$, the power exponent depends on the orbit
average, not on the chosen representative of the mantissa.

\subsection{The Lebesgue-typical law}

The doubling map is ergodic for Lebesgue measure, and
\begin{equation}
 \int_0^1\phi(t)\Differential t=-\log2=-a.                            \label{eq:phi-mean}
\end{equation}
The integral follows, for example, from the classical integral of
$\log\sin$ on $(0,\pi)$.

\begin{corollary}[Typical power exponent]
For Lebesgue-almost every $y\in(1,2)$,
\begin{equation}
 \boxed{\quad
 \log|f(2^Ny)|
 =-\frac{\bigl(\log(2^Ny)\bigr)^2}{2\log2}
 -\kappa_{\mathrm{typ}}\log(2^Ny)+\LittleO(N),
 \quad}                                                     \label{eq:typical-law}
\end{equation}
where $\kappa_{\mathrm{typ}}$ is given by \eqref{eq:ktyp}.  Equivalently,
\begin{equation}
 |f(2^Ny)|
 =\exp\!\left(-\frac{(\log(2^Ny))^2}{2\log2}\right)
 (2^Ny)^{-\kappa_{\mathrm{typ}}+\LittleO(1)}.               \label{eq:typical-multiplicative}
\end{equation}
\end{corollary}

\begin{proof}
Apply Birkhoff's ergodic theorem to $\phi\in L^1(0,1)$ and use
\eqref{eq:phi-mean} in \cref{eq:kappa-alpha}.
\end{proof}

The power in the starting draft \eqref{eq:draft-form} was
$\log(2\pi)/\log2=\log\pi/\log2+1$.  The typical exponent is larger by exactly
$1/2$.  More importantly, the error in \eqref{eq:typical-law} does not settle
to a constant or even to a deterministic logarithmic correction; it has
probabilistic fluctuations of order $\sqrt N$ and almost-sure iterated-logarithm
excursions.

\subsection{The maximal invariant average and a slow periodic ray}

The same potential appears in the classical Thue--Morse trigonometric
polynomial.  Let
\begin{equation}
 Q_N(t)=\prod_{j=0}^{N-1}2|\sin(\pi2^jt)|.
 \label{eq:QN-product}
\end{equation}
The factorization
\begin{equation}
 \prod_{j=0}^{N-1}(1-\EulerE^{2\pi\ImaginaryUnit2^jt})
 =\sum_{n=0}^{2^N-1}(-1)^{\BinaryDigitSum(n)}\EulerE^{2\pi\ImaginaryUnit nt}           \label{eq:tm-poly}
\end{equation}
shows that $Q_N$ is the modulus of a Thue--Morse polynomial.  The classical
Gelfond estimate, and its ergodic-optimization interpretation, imply that
there is a constant $C_G$ such that
\begin{equation}
 \sup_{t\in\RealNumbers}Q_N(t)\le C_G(\sqrt3)^N.                     \label{eq:gelfond}
\end{equation}
At $t=1/3$ every factor has modulus $\sqrt3$, so the exponential rate is
sharp.  In the notation for $\phi$, this says
\begin{equation}
 \sup_t\sum_{j=0}^{N-1}\phi(\T^jt)=N\beta+\BigO(1),
 \qquad \beta=\log\frac{\sqrt3}{2}.                        \label{eq:max-birkhoff}
\end{equation}
The maximizing invariant measure is supported on the period-two orbit
$\{1/3,2/3\}$; see \cite{gelfond1968,fan2021,zhang2018}.

\begin{proposition}[An exact slow ray]
Let $y=4/3$ and $x_N=2^Ny$.  Then
\begin{equation}
 |f(x_N)|
 =|f(4/3)|
 \left(\frac{3\sqrt3}{8\pi}\right)^N
 2^{-N(N+1)/2}.                                             \label{eq:exact-max-ray}
\end{equation}
Equivalently,
\begin{equation}
 |f(x_N)|=C_{4/3}
 \exp\!\left(-\frac{(\log x_N)^2}{2a}\right)x_N^{-\kappa_*},
 \label{eq:exact-max-ray-x}
\end{equation}
where
\begin{equation}
 C_{4/3}=|f(4/3)|
 \exp\!\left(\frac{(\log(4/3))^2}{2a}\right)(4/3)^{\kappa_*}>0.
\end{equation}
\end{proposition}

\begin{proof}
For every $k\ge1$,
$|\sin(\pi2^k\cdot4/3)|=\sqrt3/2$.  Substitute this into the exact finite
product preceding \eqref{eq:exact-dyadic}.  Rearranging with
$\log x_N=aN+\log(4/3)$ gives \eqref{eq:exact-max-ray-x}.
\end{proof}

\begin{remark}
The adjective ``slow'' concerns the power correction after the universal
quadratic term.  It does not mean that \eqref{eq:exact-max-ray} itself is the
largest value in the annulus $[2^N,2^{N+1}]$.  The fixed mantissa $4/3$ pays an
order-$N$ penalty from $y^{-N}$.  A larger lobe is obtained by choosing a
mantissa much closer to $1$ and allowing a transient of length about
$\log_2(N/\log N)$ before the orbit enters $\{1/3,2/3\}$.  That boundary-layer
optimization is the source of the double-log terms in \cref{sec:envelope}.
\end{remark}

\subsection{Why no global pointwise equivalent can exist}

\begin{proposition}[Failure of deterministic constant-amplitude normalization]
There is no positive function $g(x)$, finite at all sufficiently large real
$x$, for which
\begin{equation}
 g(x)|f(x)|\longrightarrow C\in(0,\infty)                  \label{eq:no-global-normalizer}
\end{equation}
as $x\to\infty$ through all real values.
Even after restricting to Lebesgue-typical dyadic rays and using the natural
centering from \eqref{eq:typical-law}, the normalized amplitude has no
almost-sure finite nonzero limit.
\end{proposition}

\begin{proof}
The first statement follows from $f(m)=0$ at every positive integer.  The
second follows from the fluctuation theorem in the next section: the centered
logarithm has arbitrarily large positive and negative excursions on almost
every ray.
\end{proof}

\section{Gaussian fluctuations and the iterated logarithm}
\label{sec:fluctuations}

The typical law gives only the order-$N$ mean.  In this particular lacunary
system, the next scale can be computed exactly.

Let $P$ be the Perron--Frobenius operator of the doubling map,
\begin{equation}
 (Ph)(x)=\frac12\left(h(x/2)+h((x+1)/2)\right).             \label{eq:transfer}
\end{equation}
A direct trigonometric calculation gives
\begin{equation}
 P\psi=\frac12\psi.                                        \label{eq:Ppsi}
\end{equation}
Define
\begin{equation}
 d=2\psi-\psi\circ\T.
 \label{eq:d-def}
\end{equation}
Using $\sin(2u)=2\sin u\cos u$,
\begin{equation}
 d(x)=\log|\tan\pi x|.                                     \label{eq:d-tan}
\end{equation}
Moreover, $Pd=0$, because the two inverse branches in \eqref{eq:transfer}
produce reciprocal tangent values.  Thus $d\circ\T^j$ is a stationary reverse
martingale-difference sequence.  From \eqref{eq:d-def},
\begin{equation}
 \sum_{j=0}^{N-1}\psi(\T^jw)
 =\sum_{j=0}^{N-1}d(\T^jw)-\psi(w)+\psi(\T^Nw).            \label{eq:martingale-coboundary}
\end{equation}
The boundary term is negligible on the $\sqrt N$ scale.

\begin{lemma}[Exact asymptotic variance]
One has
\begin{equation}
 \int_0^1d(x)^2\Differential x=\frac{\pi^2}{4}.                      \label{eq:variance}
\end{equation}
\end{lemma}

\begin{proof}
By \eqref{eq:d-tan} and symmetry,
\[
 \int_0^1d(x)^2\Differential x
 =\frac2\pi\int_0^{\pi/2}\log^2(\tan u)\Differential u.
\]
The standard beta-integral differentiation formulas give
$\int_0^{\pi/2}\log^2(\tan u)\Differential u=\pi^3/8$, yielding
\eqref{eq:variance}.  Equivalently, one may use the Fourier series
$\psi(x)=-\sum_{m\ge1}\cos(2\pi mx)/m$ and Parseval.
\end{proof}

It is useful to display the exact deterministic centering.  For
$y\in(1,2)$, put $b=\log y$ and
\begin{equation}
 C_{\mathrm{typ}}(y)
 =\log|f(y)|+\frac{b^2}{2a}+\kappa_{\mathrm{typ}}b.          \label{eq:Ctyp}
\end{equation}
Then \eqref{eq:exact-dyadic} is equivalent to
\begin{equation}
 \log|f(2^Ny)|
 +\frac{(\log(2^Ny))^2}{2a}
 +\kappa_{\mathrm{typ}}\log(2^Ny)
 -C_{\mathrm{typ}}(y)
 =\sum_{j=0}^{N-1}\psi(\T^jw).                            \label{eq:exact-centered}
\end{equation}

\begin{theorem}[Central limit theorem]
Choose $y$ uniformly from $(1,2)$.  Then, as $N\to\infty$,
\begin{equation}
 \frac{
 \log|f(2^Ny)|+
\frac{(\log(2^Ny))^2}{2a}
 +\kappa_{\mathrm{typ}}\log(2^Ny)
 }{(\pi/2)\sqrt N}
 \dist \mathcal N(0,1).                                    \label{eq:clt}
\end{equation}
\end{theorem}

\begin{proof}
The variable $w=\T(y-1)$ is Lebesgue distributed.  Since $Pd=0$,
$d\circ\T^j$ is a stationary reverse martingale-difference sequence.  It has
finite moments of every order and variance $\pi^2/4$.  The martingale central
limit theorem applies to the first sum on the right of
\eqref{eq:martingale-coboundary}; see, for example, \cite{hallheyde}.  The
coboundary term divided by $\sqrt N$ tends to zero in probability.  Finally,
$C_{\mathrm{typ}}(y)/\sqrt N\to0$ in probability; its logarithmic endpoint
singularities are square-integrable.
\end{proof}

\begin{theorem}[Law of the iterated logarithm]
For Lebesgue-almost every fixed $y\in(1,2)$, let $R_N(y)$ denote the left-hand
side of \eqref{eq:exact-centered}.  Then
\begin{align}
 \limsup_{N\to\infty}
 \frac{R_N(y)}{(\pi/2)\sqrt{2N\log\log N}}&=1,             \label{eq:lil-plus}\\
 \liminf_{N\to\infty}
 \frac{R_N(y)}{(\pi/2)\sqrt{2N\log\log N}}&=-1.           \label{eq:lil-minus}
\end{align}
\end{theorem}

\begin{proof}
The martingale law of the iterated logarithm applies to
$\sum_{j<N}d\circ\T^j$, since $d$ has moments of order $2+\delta$ for every
$\delta>0$, and the conditional quadratic variation has ergodic mean
$\pi^2/4$; see \cite{stout1970}.  The two coboundary terms in
\eqref{eq:martingale-coboundary} are negligible on the
$\sqrt{N\log\log N}$ scale.  For example, the logarithmic singularity of
$\psi$ gives exponentially small tails, so Borel--Cantelli applies.
\end{proof}

\begin{corollary}[Size of typical multiplicative oscillations]
For almost every $y$, after removal of the quadratic and typical power factors,
the remaining multiplier has subsequential logarithms of size
\begin{equation}
 \pm\frac{\pi}{2}\sqrt{2N\log\log N}.
\end{equation}
Since $N\sim\log x/\log2$, this is
\begin{equation}
 \pm\frac{\pi}{2}
 \sqrt{\frac{2\log x}{\log2}\,\log\log\log x}
 \quad\text{on the logarithmic amplitude scale}.           \label{eq:x-lil-scale}
\end{equation}
In particular, a correction of order $p\log\log x$ cannot produce a pointwise
constant normalized amplitude.
\end{corollary}

\section{The sharp dyadic-annular envelope}
\label{sec:envelope}

Define
\begin{equation}
 M_N=\max_{2^N\le x\le2^{N+1}}|f(x)|.                       \label{eq:MN}
\end{equation}
This is the largest lobe amplitude in the $N$th dyadic annulus.  The exact
identity becomes especially transparent after writing
$x=2^N(1+z)$, $0\le z\le1$:
\begin{equation}
 \begin{split}
 \log|f(2^N(1+z))|
 ={}&-\frac{aN(N+1)}2-N\log\pi\\
 &+\Bigl[\log|f(1+z)|-N\log(1+z)+S_N(z)\Bigr],             \label{eq:annular-exact}
 \end{split}
\end{equation}
where
\begin{equation}
 S_N(z)=\sum_{k=1}^{N}\phi(\T^kz).                         \label{eq:SNplus}
\end{equation}
Let
\begin{equation}
 B_N=\max_{0\le z\le1}
 \left\{\log|f(1+z)|-N\log(1+z)+S_N(z)\right\}.           \label{eq:BN}
\end{equation}
Then
\begin{equation}
 \log M_N=-\frac{aN(N+1)}2-N\log\pi+B_N.                  \label{eq:M-B}
\end{equation}

If one ignored the term $-N\log(1+z)$, the maximal Birkhoff sum would be
$N\beta+\BigO(1)$, attained on the orbit $\{1/3,2/3\}$.  But that orbit has
$z$ bounded away from zero and therefore pays a linear penalty.  The optimizer
instead starts at a tiny preimage of $1/3$.  A transient of length $r$ costs
about $ar^2/2$, while choosing $z\asymp2^{-r}$ costs about $N2^{-r}$ through
$-N\log(1+z)$.  Balancing these losses gives
$2^{-r}\asymp r/N$, hence
\begin{equation}
 r=\frac{\log N-\log\log N}{\log2}+\BigO(1).               \label{eq:r-opt}
\end{equation}

We first isolate the elementary saddle estimate.

\begin{lemma}[Boundary-layer saddle]
Fix $c>0$ and $A\in\RealNumbers$.  As $N\to\infty$,
\begin{equation}
 \min_{r\in\NaturalNumbers_0}
 \left\{\frac a2r^2+cN2^{-r}+Ar\right\}
 =\frac{(\log N)^2}{2a}
  -\frac{\log N\,\log\log N}{a}
  +\BigO(\log N).                                          \label{eq:saddle-lemma}
\end{equation}
\end{lemma}

\begin{proof}
First omit $Ar$ and minimize over real $r\ge0$.  With $q=ar$, the critical
point satisfies
\[
 q\EulerE^q=acN,
 \qquad q=W(acN),
\]
where $W$ is the principal Lambert function.  At that point the minimum is
\[
 \frac{q^2+2q}{2a}.
\]
The expansion $W(acN)=\log N-\log\log N+\BigO(1)$ gives the two displayed
terms and an error $\BigO(\log N)$.  Rounding $r$ to an integer changes the
value by $\BigO(\log N)$, because the second derivative at the saddle is
$\BigO(\log N)$.  Any minimizing integer is $\BigO(\log N)$; the added term
$Ar$ therefore also changes the result by only $\BigO(\log N)$.
\end{proof}

\begin{theorem}[Sharp dyadic-annular maximum]
As $N\to\infty$,
\begin{equation}
 \boxed{\begin{aligned}
 \log M_N={}&-\frac a2N^2
 +\bigl(\beta-\log\pi-a/2\bigr)N\\
 &-\frac{(\log N)^2}{2a}
 +\frac{\log N\,\log\log N}{a}
 +\BigO(\log N).
 \end{aligned}}                                             \label{eq:MN-sharp}
\end{equation}
\end{theorem}

\begin{proof}
We prove matching bounds for $B_N$.

\emph{Upper bound.}
The product representation gives $|f(1+z)|\le1$.  If $z>1/4$, then
$-N\log(1+z)\le-c_0N$ for an absolute $c_0>0$, while
\eqref{eq:max-birkhoff} gives $S_N(z)\le N\beta+\BigO(1)$.
This region is far below the eventual maximum, whose loss from $N\beta$ is
only of order $(\log N)^2$.

Now suppose $0<z\le1/4$, and let $r$ be the largest nonnegative integer such
that
\begin{equation}
 2^rz\le\frac14.                                            \label{eq:r-def}
\end{equation}
If $r\ge N$, then all $N$ factors lie in the small-angle regime and
$S_N(z)\le-aN^2/2+\BigO(N)$, again much too small.  We may therefore assume
$r<N$.  For $1\le k\le r$ there is no wrap modulo one, and
$\sin(\pi2^kz)\le\pi2^kz$.  Because
$z\le2^{-r-2}$,
\begin{align}
 \sum_{k=1}^{r}\phi(\T^kz)
 &\le r\log\pi+r\log z+\frac{ar(r+1)}2\notag\\
 &\le-\frac a2r^2+r\left(\log\pi-\frac{3a}{2}\right).     \label{eq:transient-upper}
\end{align}
The remaining $N-r$ terms satisfy the shifted Gelfond estimate
\[
 \sum_{k=r+1}^{N}\phi(\T^kz)
 \le(N-r)\beta+\BigO(1).
\]
Maximality in \eqref{eq:r-def} also gives $z>2^{-r-3}$.  Since
$\log(1+z)\ge z/2$ for $0\le z\le1$,
\[
 N\log(1+z)\ge2^{-r-4}N.
\]
Combining the last three estimates, for fixed constants $c,C>0$,
\begin{equation}
 \log|f(1+z)|-N\log(1+z)+S_N(z)
 \le N\beta-\frac a2r^2-cN2^{-r}+Cr+C.                    \label{eq:BN-upper-r}
\end{equation}
The saddle lemma implies
\begin{equation}
 B_N\le N\beta-
\frac{(\log N)^2}{2a}
 +\frac{\log N\,\log\log N}{a}+\BigO(\log N).             \label{eq:BN-upper}
\end{equation}

\emph{Lower bound.}
For an integer $1\le r<N$, take the explicit preperiodic point
\begin{equation}
 z_r=\frac{1}{3\,2^r}.                                      \label{eq:zr}
\end{equation}
After $r$ iterations it reaches $1/3$ and then alternates between $1/3$ and
$2/3$.  Therefore
\begin{equation}
 S_N(z_r)=A_r+(N-r)\beta,                                   \label{eq:SN-zr}
\end{equation}
where
\begin{equation}
 A_r=\sum_{j=0}^{r-1}\log\sin\!\left(\frac{\pi}{3\,2^j}\right).
\end{equation}
Separating each small angle from its sinc correction gives
\begin{align}
 A_r
 & =r\log(\pi/3)-\frac{ar(r-1)}2
 +\sum_{j=0}^{r-1}\log\SincRad\!\left(\frac{\pi}{3\,2^j}\right)\notag\\
 & =r\log(\pi/3)-\frac{ar(r-1)}2+\log f(1/3)+\LittleO(1).   \label{eq:Ar-asymp}
\end{align}
The zero of $f$ at $1$ is simple, so
\begin{equation}
 \log|f(1+z_r)|=\log z_r+\BigO(1)=-ar+\BigO(1).             \label{eq:f-near-one}
\end{equation}
Also, $\log(1+z_r)\le z_r$.  Substitution in \eqref{eq:BN} yields
\begin{equation}
 B_N\ge N\beta-\frac a2r^2-\frac{N}{3\,2^r}+\BigO(r).     \label{eq:BN-lower-r}
\end{equation}
Choose $r$ according to the saddle lemma.  Then
\begin{equation}
 B_N\ge N\beta-
\frac{(\log N)^2}{2a}
 +\frac{\log N\,\log\log N}{a}+\BigO(\log N).             \label{eq:BN-lower}
\end{equation}
Together, \eqref{eq:BN-upper} and \eqref{eq:BN-lower} determine $B_N$ to the
required precision.  Finally use \eqref{eq:M-B}.
\end{proof}

\subsection{The envelope in the original variable}

Let $X=2^N$ and $L=\log X$.  Since $N=L/a$, \eqref{eq:MN-sharp} becomes
\begin{equation}
 \boxed{\begin{aligned}
 \log\max_{X\le x\le2X}|f(x)|
 ={}&-\frac{L^2}{2a}-\kappa_*L\\
 &-\frac{(\log L)^2}{2a}
 +\frac{\log L\,\log\log L}{a}
 +\BigO(\log L),
 \end{aligned}}                                             \label{eq:MX-sharp}
\end{equation}
for dyadic $X\to\infty$.  Equivalently,
\begin{equation}
 \begin{split}
 \max_{X\le x\le2X}|f(x)|
 ={}&\exp\!\left(-\frac{(\log X)^2}{2\log2}\right)
 X^{-\kappa_*}\\
 &\times\exp\!\left(
 -\frac{(\log\log X)^2}{2\log2}
 +\frac{\log\log X\,\log\log\log X}{\log2}
 +\BigO(\log\log X)\right).
 \end{split}                                                \label{eq:MX-multiplicative}
\end{equation}
The final exponential factor is smaller than every fixed negative power of
$\log X$, but larger than every fixed negative power of $X$.  It is therefore
a genuinely intermediate correction, invisible in an ansatz containing only
$(\log X)^{-p}$.

\begin{corollary}[Near-optimal boundary layer]
The explicit points used in the lower bound have
\begin{equation}
 x_{N,r}=2^N\left(1+\frac{1}{3\,2^r}\right),
 \qquad
 r=\log_2N-\log_2\log N+\BigO(1).                          \label{eq:near-opt-location}
\end{equation}
Consequently their relative displacement from the left endpoint is
\begin{equation}
 \frac{x_{N,r}-2^N}{2^N}
 =\Theta\!\left(\frac{\log N}{N}\right)
 =\Theta\!\left(\frac{\log\log X}{\log X}\right).         \label{eq:relative-location}
\end{equation}
\end{corollary}

\begin{proof}
Equation \eqref{eq:r-opt} gives $2^{-r}=\Theta((\log N)/N)$, and
\eqref{eq:relative-location} follows from \eqref{eq:zr}.
\end{proof}

The theorem does not assert that every exact maximizer is one of the special
preimages \eqref{eq:zr}; rather, those points construct the correct lower
asymptotic, while the upper proof shows that any competing point must pay the
same boundary-layer cost to the displayed precision.

\section{Comparison of the distinct decay regimes}

The formulas are easiest to compare after extracting the universal factor
\begin{equation}
 \mathcal L(x)=\exp\!\left(-\frac{(\log x)^2}{2\log2}\right).
\end{equation}

\begin{center}
\begin{tabularx}{0.98\textwidth}{@{}>{\raggedright\arraybackslash}p{0.22\textwidth}
 >{\raggedright\arraybackslash}p{0.26\textwidth}X@{}}
\toprule
Regime & Power after $\mathcal L(x)$ & Additional behavior\\
\midrule
Lebesgue-typical dyadic ray
& $x^{-\kappa_{\mathrm{typ}}}$, $\kappa_{\mathrm{typ}}\approx3.151496$
& Log-amplitude fluctuations have standard deviation
  $(\pi/2)\sqrt{\log_2x}$ and obey an iterated-logarithm law.\\[0.5em]
Maximal invariant periodic ray
& $x^{-\kappa_*}$, $\kappa_*\approx2.359015$
& The ray $x=(4/3)2^N$ has an exact constant multiplier after this normalization.\\[0.5em]
Largest lobe in $[X,2X]$
& $X^{-\kappa_*}$
& Multiplied by the boundary-layer factor in \eqref{eq:MX-multiplicative}, whose leading logarithm is
  $-(\log\log X)^2/(2\log2)$.\\[0.5em]
Starting-draft ansatz
& $x^{-\log(2\pi)/\log2}$, exponent $\approx2.651496$
& A fixed power of $\log x$ cannot match either typical fluctuations or the annular boundary layer.\\
\bottomrule
\end{tabularx}
\end{center}

Several conclusions follow.

\begin{enumerate}
\item The coefficient $1/(2\log2)$ of $(\log x)^2$ is robust in all regular
      regimes because it comes from the deterministic denominator
      $\prod_{k\le N}2^k$.
\item The coefficient of $\log x$ is not universal.  It is the affine image
      \eqref{eq:kappa-alpha} of a Birkhoff average.
\item At typical points, the next visible term is not a deterministic
      $\log\log x$ term.  Random-looking lacunary fluctuations are much larger.
\item For the largest lobe, optimizing the time needed to enter the maximizing
      period-two orbit creates a quadratic function of $\log\log x$.
\item Because of exact integer zeros and arbitrarily deep near-zero returns,
      there is no nontrivial global lower envelope comparable to the upper one.
\end{enumerate}

\section{A general geometric ratio}

The mechanism is not special to the notation of the up-function.  For an
integer $q\ge2$, consider
\begin{equation}
 f_q(x)=\prod_{n=0}^{\infty}
 \SincRad\!\left(\frac{\pi x}{q^n}\right).                    \label{eq:fq}
\end{equation}
Writing $x=q^Ny$ gives the exact analogue
\begin{equation}
 \log|f_q(q^Ny)|
 =\log|f_q(y)|-N\log(\pi y)
 -\frac{\log q}{2}N(N+1)
 +\sum_{k=1}^{N}\phi(T_q^kz),                              \label{eq:q-decomp}
\end{equation}
where $T_qt=qt\pmod1$ and $z$ is the fractional part of $y$.  Hence the
universal quadratic coefficient becomes
\begin{equation}
 -\frac{(\log x)^2}{2\log q}.                              \label{eq:q-leading}
\end{equation}
For Lebesgue-typical mantissas the average of $\phi$ is still $-\log2$, so the
typical power exponent is
\begin{equation}
 \kappa_{\mathrm{typ}}(q)
 =\frac{\log\pi+(\log q)/2+\log2}{\log q}.                 \label{eq:q-typ}
\end{equation}
The slow envelope is controlled by the maximal invariant average of $\phi$
under $T_q$.  For $q=2$, the exceptional simplification is that the exact
maximum is $\log(\sqrt3/2)$ and is realized by the period-two orbit
$\{1/3,2/3\}$, linking the problem directly to the classical Thue--Morse
Gelfond exponent.

\section{Conclusions}

The Fourier transform of Rvachev's up-function has a clean leading scale but no
single clean pointwise equivalent.  The exact finite decomposition
\eqref{eq:exact-dyadic} is the organizing principle:
\begin{equation*}
 \text{quadratic denominator cost}
 \quad+\quad
 \text{doubling-map Birkhoff sum}.
\end{equation*}
It yields the following rigorous answer to the decay question.

\begin{itemize}
\item Uniformly, $|f(x)|$ is at most
      $\exp(-(\log x)^2/(2\log2)+\BigO(\log x))$, hence it is
      super-polynomially small.
\item On almost every fixed dyadic ray, the power exponent after this
      log-normal factor is $\kappa_{\mathrm{typ}}$, but the logarithm retains
      Gaussian and iterated-logarithm fluctuations.
\item The slowest invariant average is the Thue--Morse period-two average,
      giving exponent $\kappa_*$.  It is realized exactly on rays such as
      $(4/3)2^N$.
\item The largest lobe in a dyadic annulus has the sharper expansion
      \eqref{eq:MN-sharp}--\eqref{eq:MX-multiplicative}; its correction is
      quadratic in $\log\log x$, not a fixed power of $\log x$.
\end{itemize}

Thus the proposed constant-amplitude normalization in the starting draft must
be replaced by a regime-dependent statement.  The typical law, exceptional
periodic law, fluctuation theorem, and annular envelope together give a precise
and internally consistent description of how fast the infinite sinc product
decays.

\appendix

\section{Auxiliary integral and covariance checks}

For completeness, the Fourier series
\begin{equation}
 \psi(x)=\log|2\sin\pi x|
 =-\sum_{m=1}^{\infty}\frac{\cos(2\pi mx)}{m}              \label{eq:psi-fourier}
\end{equation}
holds in $L^2(0,1)$.  Parseval gives
\begin{equation}
 \int_0^1\psi(x)^2\Differential x=\frac12\sum_{m=1}^{\infty}\frac1{m^2}
 =\frac{\pi^2}{12}.                                        \label{eq:psi-var}
\end{equation}
Because only frequencies divisible by $2^k$ survive when
$\psi$ is paired with $\psi\circ\T^k$,
\begin{equation}
 \int_0^1\psi(x)\psi(\T^kx)\Differential x
 =\frac{\pi^2}{12\,2^k}.                                  \label{eq:covariance}
\end{equation}
Therefore the Green--Kubo variance is
\begin{equation}
 \frac{\pi^2}{12}
 +2\sum_{k=1}^{\infty}\frac{\pi^2}{12\,2^k}
 =\frac{\pi^2}{4},                                         \label{eq:green-kubo}
\end{equation}
in agreement with \eqref{eq:variance} and the martingale decomposition.

\section{A direct formula for the simple zero at one}

From \eqref{eq:weier-product},
\begin{equation}
 f(x)=(1-x^2)G(x),
 \qquad
 G(x)=\prod_{m=2}^{\infty}
 \left(1-\frac{x^2}{m^2}\right)^{1+\TwoAdicValuation(m)}.              \label{eq:G-near-one}
\end{equation}
The product for $G$ converges locally uniformly near $x=1$ and $G(1)\ne0$.
Consequently,
\begin{equation}
 f(1+z)=-2G(1)z+\BigO(z^2),                                \label{eq:simple-zero}
\end{equation}
which justifies \eqref{eq:f-near-one} and quantifies the endpoint cost in the
annular lower bound.

\begin{thebibliography}{99}

\bibitem{arias1982}
J.~Arias de Reyna,
\newblock \emph{An infinitely differentiable function with compact support:
Definition and properties},
\newblock Rev. Real Acad. Cienc. Exactas Fis. Natur. Madrid \textbf{76} (1982),
21--38; English translation, arXiv:1702.05442 (2017),
\newblock \url{https://arxiv.org/abs/1702.05442}.

\bibitem{fan2021}
A.~Fan, J.~Schmeling, and W.~Shen,
\newblock \emph{$L^\infty$-estimation of generalized Thue--Morse
trigonometric polynomials and ergodic maximization},
\newblock Discrete Contin. Dyn. Syst. \textbf{41} (2021), 297--327,
\newblock \url{https://doi.org/10.3934/dcds.2020363}.

\bibitem{gelfond1968}
A.~O. Gel'fond,
\newblock \emph{Sur les nombres qui ont des proprietes additives et
multiplicatives donnees},
\newblock Acta Arith. \textbf{13} (1967/68), 259--265,
\newblock \url{https://doi.org/10.4064/aa-13-3-259-265}.

\bibitem{hallheyde}
P.~Hall and C.~C. Heyde,
\newblock \emph{Martingale Limit Theory and Its Application},
\newblock Academic Press, New York, 1980.

\bibitem{proveit}
V.~Reshetnikov,
\newblock \emph{ProveIt: Analysis/FabiusFunction}, GitHub repository,
\newblock \url{https://github.com/VladimirReshetnikov/ProveIt/tree/main/Analysis/FabiusFunction}.

\bibitem{decaydraft}
V.~Reshetnikov,
\newblock \emph{Decay rate conjecture}, research-frontier draft in the
\texttt{ProveIt} repository,
\newblock \url{https://github.com/VladimirReshetnikov/ProveIt/blob/main/Analysis/FabiusFunction/docs/semi-formalized-research-frontiers/drafts/Decay%20rate%20conjecture/Decay%20rate%20conjecture.tex}.

\bibitem{stout1970}
W.~F. Stout,
\newblock \emph{A martingale analogue of Kolmogorov's law of the iterated
logarithm},
\newblock Z. Wahrscheinlichkeitstheorie verw. Gebiete \textbf{15} (1970),
279--290.

\bibitem{wikipedia}
Wikipedia contributors,
\newblock \emph{Fabius function}, section ``Fourier transform'',
\newblock \url{https://en.wikipedia.org/wiki/Fabius_function}.

\bibitem{zhang2018}
Y.~Zhang, K.~Yin, and W.~Wu,
\newblock \emph{A rigorous computer aided estimation for Gelfond exponent of
weighted Thue--Morse sequences},
\newblock arXiv:1806.08329 (2018),
\newblock \url{https://arxiv.org/abs/1806.08329}.

\end{thebibliography}

\end{document}
